\chapter{从点云几何到相对位姿：GICP}

\label{chap:md-11}

机器人在两个位置各采集一帧点云。即使两帧看到了同一堵墙，点的坐标也不会自动相同：传感器的位置和朝向变了，每个点还分别在当时的传感器坐标系中表达。要把两帧放到一起，需要找一个位姿变换；若这个变换合理，同一处几何结构应当互相贴合。GICP（Generalized Iterative Closest Point）就是把这种贴合写成一个可以迭代求解的局部优化问题。

上一章的预积分从一串高频IMU读数得到区间相对运动；本章从两份点云得到几何相对位姿。两者都可能被用于关键帧之间的约束，但来源、误差结构和失效方式不同。本章沿“点在哪个坐标系—如何找对应—怎样度量对应误差—怎样在$SE(3)$上更新—结果何时可信”走一遍，不把ICP当作一个只需调用的黑箱。GICP的局部协方差与目标形式可参见\href{https://www.roboticsproceedings.org/rss05/p21.html}{Segal等人的原论文}。

\section{先把两帧点云放在同一个问题中}

\subsection{源点、目标点和待估计变换}

记源点云为$\mathcal P=\{p_k\}$，点$p_k$在源坐标系$S$中表达；目标点云为$\mathcal Q=\{q_l\}$，点$q_l$在目标坐标系$Q$中表达。本章求$T_{QS}=(R,t)\in SE(3)$，它把源系坐标送到目标系：

\begin{equation*}
y_k=T_{QS}\cdot p_k=Rp_k+t.
\end{equation*}

这里的$Q$是目标坐标系的名称，不是噪声协方差。若把源和目标交换，所求位姿也要取逆，残差与Jacobian的写法随之改变。先固定方向，比记住某个库的参数顺序可靠。

在关键帧图中，常用$T_i=T_{Wi}$表示第$i$帧到世界系的变换。于是从第$j$帧到第$i$帧的预测相对变换为$T_i^{-1}T_j$。若GICP把$j$帧作为源、$i$帧作为目标，它得到的测量才可直接与这个预测比较。目标若是局部子地图，还要说清子地图使用哪一个坐标系。

\subsection{一对点为什么还不够}

假设$p_k$和$q_k$确实来自环境中的同一处。直观上希望$Rp_k+t$靠近$q_k$，可以先写差

\begin{equation*}
e_k(T)=Rp_k+t-q_k\in\mathbb R^3.
\end{equation*}

一个点对只约束变换后的一个位置；旋转如何确定，要依靠不同空间位置的多对点。若大多数点都落在一张近似平面上，某些平移和转动方向仍可能很弱。后面将用局部协方差和近似Hessian看见这种差别。

\subsection{对应关系本身也未知}

实际数据没有预先标好的$p_k\leftrightarrow q_k$。给定当前变换$T$，先把$p_k$变到目标系，再在目标点云或目标地图中找候选邻居；求出新的位姿后，变换后的点位置改变，对应关系也可能改变。因此配准通常交替做：

\begin{enumerate}

\item 用当前$T$寻找对应；

\item 在暂时固定的对应下减小残差；

\item 更新$T$并重新寻找对应。

\end{enumerate}

最近邻不保证是真正的同一处。稀疏采样、重复墙面、动态物体和较差初值都可能给出错误对应。

\section{从点到点、点到平面走向GICP}

\subsection{点到点：三个坐标方向一视同仁}

最简单的目标把每对点的欧氏差平方后相加：

\begin{equation*}
E_{\mathrm{p2p}}(T)=\sum_k\|Rp_k+t-q_k\|^2.
\end{equation*}

它把目标点看作必须在三个方向都靠近的一个位置。真实LiDAR采到同一片墙时，两帧的离散采样通常并不落在完全相同的墙面位置；沿墙面的差可以较大，却不一定表示位姿错误。

\subsection{点到平面：先看法线方向}

若目标点附近是一片平面，估计单位法向$n_k$，可以只比较法向距离：

\begin{equation*}
r_k^{\mathrm{plane}}=n_k^\top(Rp_k+t-q_k).
\end{equation*}

沿法线离开墙面会产生明显误差；沿墙面滑动时，这个残差可能很小。这比逐坐标比较更接近“两个点是否落在同一片局部表面”的问题，但也说明只看到一张墙时，沿墙面的位移很难由该墙单独确定。

\subsection{局部协方差：把允许偏离的方向写出来}

对$p_k$、$q_k$各取一个局部邻域。邻域点的散布可以给出局部主方向：近似平面通常在切向分散较大，在法向分散较小；边缘或角点的形状又不同。记源侧和目标侧的正定、经过适当正则化的局部协方差为$C_{p_k}$与$C_{q_k}$。它们描述配准模型允许点在各方向偏离多少，不能直接等同于LiDAR单次测距的传感器噪声协方差。

\begin{bookexample}{一面水平地面}

如果邻域在水平面内延伸，而竖直方向很薄，协方差的两个水平特征值较大，竖直特征值较小。两帧点沿地面相差一些距离，惩罚较轻；其中一帧整体抬高，竖直残差会被更明显地惩罚。这个例子只解释一对局部表面的权重，不能据此断言整帧六个自由度都已可观。

\end{bookexample}

\section{GICP怎样给一个点对加权}

\subsection{把两侧局部不确定性写到目标坐标系}

源侧协方差原本在源系。旋转$R$后，它在目标系中的表示是$RC_{p_k}R^\top$；目标侧协方差已经在目标系。采用两侧局部偏离独立的模型时，差$e_k$对应的协方差为

\begin{equation*}
\Sigma_k(T)=C_{q_k}+RC_{p_k}R^\top.
\end{equation*}

平移不会改变协方差的方向，旋转会。因此$e_k$和$\Sigma_k$都依赖待求位姿，虽然依赖方式不同。若局部邻域给出的矩阵半正定或病态，求逆之前需要正则化；直接计算显式逆在数值实现中也不是首选。

\subsection{同样大小的残差为什么受到不同惩罚}

局部高斯模型下，一个对应的常用加权误差为

\begin{equation*}
E_k(T)=e_k(T)^\top\Sigma_k(T)^{-1}e_k(T).
\end{equation*}

把$\Sigma_k=L_kL_k^\top$分解后，令$L_k\widetilde e_k=e_k$，就得到$E_k=\|\widetilde e_k\|^2$。这是第\ref{chap:md-10}章白化思想在点云中的另一种使用：先说明误差在各方向可信到什么程度，再比较变换后的点。

需要分清目标函数与严格概率密度。若把状态相关的$\Sigma_k(T)$当作完整高斯协方差，负对数还含$\log\det\Sigma_k(T)$项；常见GICP配准目标使用上述Mahalanobis项，并在迭代中按具体实现处理协方差。这里讨论的是这一配准目标的局部优化，不能无条件把它称为所有原始点云观测的精确似然。

\subsection{多个对应与鲁棒处理}

将合格点对集合记为$\mathcal C$，可写

\begin{equation*}
E(T)=\sum_{k\in\mathcal C}\rho\!\left(e_k^\top\Sigma_k^{-1}e_k\right),
\end{equation*}

其中$\rho$是所选鲁棒代价。对应过滤、最大匹配距离、法向一致性、动态点剔除和鲁棒核分别处理不同来源的错误；它们不能使错误的地点匹配自动变成正确约束。有效点对变少时，还需检查剩余几何是否足以支撑六维位姿。

\section{在$SE(3)$上怎样走出一次优化迭代}

\subsection{先固定局部扰动约定}

本节将$\delta\xi=(\delta\rho^\top,\delta\phi^\top)^\top$按平移在前、旋转在后排列，采用左侧扰动

\begin{equation*}
T'=\operatorname{Exp}(\delta\xi^\wedge)T.
\end{equation*}

这与第\ref{chap:md-06}章的左侧局部参数化一致。它是本节的计算约定；第八、九章的惯性状态姿态使用右侧误差。两个模块可以采用不同局部坐标，但把它们连接时必须做相应换元，不能把Jacobian直接按数组位置拼接。

\subsection{一个点对的残差怎样变化}

在当前$T$下令$y_k=Rp_k+t$。对很小的$\delta\xi$，有

\begin{equation*}
\operatorname{Exp}(\delta\xi^\wedge)\cdot y_k
\simeq y_k+\delta\rho+\delta\phi\times y_k
=y_k+\delta\rho-y_k^\wedge\delta\phi.
\end{equation*}

在本次线性化中暂时固定对应$q_k$，点差的一阶变化为

\begin{equation*}
e_k(T')\simeq e_k(T)+J_k\delta\xi,\qquad
J_k=\begin{bmatrix}I&-y_k^\wedge\end{bmatrix}.
\end{equation*}

这张$3\times6$矩阵把位姿附近的六维局部增量送到目标系三维点差。它是点差的Jacobian；若在同一次微分中还保留$\Sigma_k(T)$对$R$的变化，目标函数的导数另有相应项。若选右侧扰动，$J_k$也要重新求，不能只改变最后的位姿更新式。

\subsection{多个点对怎样形成局部方程}

暂时固定对应和当前权重，线性化后的加权二次模型汇总为

\begin{equation*}
H\simeq\sum_kJ_k^\top W_kJ_k,\qquad
g\simeq\sum_kJ_k^\top W_ke_k,\qquad
H\delta\xi=-g,
\end{equation*}

其中$W_k$由$\Sigma_k^{-1}$及鲁棒权重构成。解出$\delta\xi$后，用指数映射得到合法的$T'$，再重新找对应和评价目标。求解时通常使用矩阵分解或阻尼，而不是显式计算$H^{-1}$。上式是固定对应与权重的一步局部模型，不保证从任意初值找到全局最优。

\subsection{初值为什么会改变结果}

若初始$T$已经使同一处墙面相互靠近，最近邻更可能找对，线性化也更可能处在合理区域。若初值把源点放到相似的另一条走廊，算法可能从错误对应出发并收敛到错误的局部解。IMU预测、上一帧里程计或地点识别给出的粗位姿，常用于提供初值；这不等于它们自动成为GICP输出之外的一份独立观测。

\section{怎样判断配准结果是否值得使用}

\subsection{残差小不一定代表六个方向都确定}

假设只看到一面无限延伸的平墙。沿墙面的平移可能几乎不改变点到平面的误差，因此即使总误差很小，相应方向仍可能没有得到足够约束。局部近似Hessian的很小特征值提示弱约束方向；重复结构和错误对应也可能造成“看起来收敛”的假象。检查时应同时看残差、有效对应数量、重叠率、位姿增量和几何退化。

\subsection{从局部曲率到位姿不确定性需要条件}

在对应关系可信、噪声权重合理、线性化误差较小且$H$非奇异时，$H^{-1}$可以作为局部位姿协方差的一个近似。它没有自动包含错误对应、动态物体、初值选择和模型误差；退化时更不能把病态矩阵硬求逆后当作可靠权重交给后端。系统可以拒绝本次配准、减弱弱方向的权重，或依靠其他传感器补充约束。

\subsection{点云与子地图配准}

当前扫描可以对上一帧配准，也可以对由多个关键帧组成的局部子地图配准。后一种方式通常提供更丰富的几何，但目标地图会随关键帧位姿和回环校正而变化。要保存子地图坐标系、参与构图的关键帧及其版本，否则后端调整后，前端可能还在使用旧几何。

\section{GICP在系统里交出什么}

\subsection{作为实时状态修正}

若配准得到当前扫描相对于局部地图的位姿，前端可以把它作为第\ref{chap:md-08}章的外部观测，再用所选观测模型和协方差做ESKF更新。另一种设计直接用点残差更新状态；此时状态Jacobian、对应关系和迭代次序都要在滤波器里处理。这里仅说明两种接口，不把同一组点同时当作两份独立观测。

\subsection{作为关键帧之间的几何因子}

若第$i,j$个关键帧位姿为$T_i,T_j$，GICP测得$j$到$i$的相对变换$Z_{ij}$，则可以定义

\begin{equation*}
r_{ij}^{\mathrm{geom}}
=\operatorname{Log}\!\left(Z_{ij}^{-1}T_i^{-1}T_j\right)^\vee.
\end{equation*}

这个六维残差比较“点云给出的相对位姿”与“当前两端状态给出的相对位姿”。它的协方差必须在相同的局部残差坐标中表达。当前配准结果也可以只用来初始化后端；是否另立一个因子，要结合数据相关性与系统设计决定。

\subsection{作为回环的几何验证}

地点识别可以给出“当前地点可能是历史地点”的候选，但尚未给出可信的六维约束。用两个关键帧或子地图运行GICP、检查重叠和退化，才可能生成回环相对位姿。错误回环会把时间上相隔很远的状态错误地拉到一起，因此验证标准应比普通相邻帧配准更严格。完整流程见第\ref{chap:md-12}章。

\section{回到这两帧点云}

\begin{center}\small\begin{tabular}{p{0.420\textwidth}p{0.420\textwidth}}\toprule

问题 & 本章的答案 \\

\midrule

输入是什么 & 两帧点云或扫描与子地图，以及一个初始相对位姿 \\

求的是什么 & 把源系送到目标系的$T_{QS}\in SE(3)$ \\

残差在哪里 & 变换后点与目标点的三维差，或其局部几何投影 \\

权重来自哪里 & 点邻域的局部协方差、对应筛选与鲁棒模型 \\

何时可用于后端 & 对应、重叠、残差和弱约束方向经过检查后 \\

\bottomrule\end{tabular}\end{center}

GICP给出的是由当前点云和目标几何支持的局部相对位姿。IMU预积分给出的是由一段惯性读数支持的运动约束。下一章将追踪这两份信息怎样进入同一个SLAM系统，以及回环发生时历史状态与地图怎样一起改变。
